% TOP_PROBLEM: {"id": 2522, "title": "Kourovka Notebook Problem 21.13", "category_id": 17, "category": {"id": 17, "name": "group_theory", "display_name": "Group Theory", "description": "Problems about groups, group actions, representations, and related algebraic structures.", "slug": "group-theory", "order_index": 17, "created_at": "2026-07-31T15:26:25.671Z"}, "status": "open", "problem_number": "KOU-21.13", "set_id": 10, "statusQualification": "The relative phrase still larger is resolved using the continuum-rank benchmark stated in the original background."}
% FIRST_POSTED: 2026-10-01
% ENTRY_KIND: statement_only
% UnsolvedMath ID 2522; source catalogue code KOU-21.13
% !TeX program = lualatex
% Definition and sources only; this file is not an LLM solution attempt.
\documentclass[11pt,a4paper]{article}
\usepackage[margin=25mm]{geometry}
\usepackage{fontspec}
\setmainfont{lmroman10-regular.otf}[BoldFont=lmroman10-bold.otf,
 ItalicFont=lmroman10-italic.otf,BoldItalicFont=lmroman10-bolditalic.otf]
\usepackage{amsmath,amssymb,mathtools}
\usepackage{unicode-math}
\setmathfont{latinmodern-math.otf}
\AtBeginDocument{\let\setminus\smallsetminus}
\usepackage{xurl,hyperref}
\hypersetup{colorlinks=true,urlcolor=blue}
\setcounter{secnumdepth}{0}
\setlength{\parindent}{0pt}
\setlength{\parskip}{5pt}
\setlength{\emergencystretch}{3em}
\begin{document}
\section*{Kourovka Notebook Problem 21.13}\label{2522}
{\small UnsolvedMath ID 2522 · KOU-21.13 · Group Theory · Kourovka Notebook - New Problems, Issue 21}
\subsection{Definitions and mathematical statement}
Let $\omega=\{0,1,2,\ldots\}$ and let
$\mathbb Z^\omega=\prod_{i\in\omega}\mathbb Z$ be the unrestricted
product, with coordinatewise addition. Its elements are all integer
sequences, without a finite-support condition. For an abelian subgroup
$A\le\mathbb Z^\omega$, define its algebraic dual by
\[
 A^*=\operatorname{Hom}_{\mathbb Z}(A,\mathbb Z),
 \qquad (f+g)(a)=f(a)+g(a).
\]
Homomorphisms are abstract additive maps. In particular this is the
integer-valued algebraic dual, not a dual defined using the circle
group or a topology.

Write $\mathfrak c=2^{\aleph_0}$ for the cardinality of the continuum.
Does there exist a subgroup $A\le\mathbb Z^\omega$ for which $A^*$ is
free abelian of rank $\kappa>\mathfrak c$? Freeness means that some
family $(f_\alpha)_{\alpha<\kappa}$ is a basis: every member of $A^*$
has a unique expression as a finite integer linear combination of
these homomorphisms. Equivalently,
$A^*\cong\bigoplus_{\alpha<\kappa}\mathbb Z$.

The known continuum-rank case motivates the strict inequality; it
would not answer the question. The subgroup $A$ itself is not required
to be free, countable, or of finite index. Since $|A|\le\mathfrak c$,
there are at most $\aleph_0^{\mathfrak c}=2^{\mathfrak c}$ maps
$A\to\mathbb Z$, giving the upper bound
$\kappa\le2^{\mathfrak c}$. Attaining that upper bound is not required:
any free dual rank strictly above $\mathfrak c$ would suffice.
\subsection{Short English statement}
Can a subgroup of the countable product of integers have an integer-valued algebraic dual that is free abelian of rank above the continuum?
\subsection{Sources}
\begin{itemize}
\item[S1] ULAM.AI and the UnsolvedMath Contributors, supplied problems.json, record 2522 (KOU-21.13), Kourovka Notebook - New Problems, Issue 21. Catalogue identity and source transcription; CC BY 4.0.
\url{https://huggingface.co/datasets/ulamai/UnsolvedMath}.
\item[S2] The Kourovka Notebook, 21st edition, new problems, Problem 21.13. Expanded formulation of the supplied catalogue statement.
\url{https://arxiv.org/pdf/1401.0300v47}.
\end{itemize}
\end{document}
